
\vspace{0cm}

% ── A1. Introducción ────────────────────────────────────────────────
\Large{\textsc{\textcolor{nar}{A1. Introducción}}}
\vspace{0.2cm}

\normalsize{Este Trabajo Fin de Grado desarrolla y optimiza Voctomix~2.0, un sistema de código
abierto para la producción remota de vídeo en tiempo real. El proyecto amplía la herramienta
Voctomix con nuevas funcionalidades: composición de vídeo multicámara, \textit{overlays} dinámicos,
telemetría basada en RabbitMQ y despliegue con Docker y Kubernetes, con el objetivo de que
cualquier organización pueda realizar retransmisiones de calidad profesional sin necesidad de
\textit{hardware} dedicado.

El sistema nació de una necesidad real: muchos grupos universitarios, asociaciones y entidades
públicas necesitan cubrir eventos en directo pero no pueden permitirse el equipamiento \textit{broadcast}
habitual. Al estar construido sobre herramientas libres, como GStreamer, Python, Docker, RabbitMQ
y Kubernetes, el proyecto es auditable, reproducible y puede evolucionar más allá del trabajo
individual, alineándose con los principios del \textit{software} libre y la ingeniería abierta.

Este anexo revisa los aspectos éticos, económicos, sociales y ambientales del proyecto.}
\vspace{0.3cm}

% ── A2. Impactos relevantes ──────────────────────────────────────────
\Large{\textsc{\textcolor{nar}{A2. Descripción de impactos relevantes relacionados con el proyecto}}}
\vspace{0.2cm}

\normalsize{A continuación se identifican y justifican los principales impactos del proyecto,
agrupados por dimensión:}

\vspace{0.15cm}
\noindent\textbf{Impacto en el sector audiovisual y académico}\\
\normalsize{Voctomix~2.0 está pensado para organizaciones que necesitan producción audiovisual pero
no tienen presupuesto \textit{broadcast}: grupos de investigación universitarios como CyberNEMO (UPM),
asociaciones estudiantiles, organismos públicos y entidades sin ánimo de lucro. Con él pueden
hacer retransmisiones multicámara con transiciones animadas, \textit{overlays}, rotulación dinámica y
grabación simultánea, usando equipos convencionales, sin \textit{hardware} dedicado.}

\vspace{0.15cm}
\noindent\textbf{Aspectos éticos}\\
\normalsize{El sistema está construido sobre \textit{software} libre: Voctomix tiene licencia MIT,
GStreamer es LGPL~2.1 y Docker y Kubernetes usan Apache~2.0~\cite{apache2_license}.
Que todo el código sea abierto y auditable significa que no hay dependencia de ningún
proveedor propietario.

Dicho esto, el modelo de producción remota tiene implicaciones que conviene señalar. Al poder
gestionar toda la producción desde un equipo centralizado, se reduce la necesidad de personal
técnico presente en el lugar del evento: operadores de cámara, técnicos de audio o encargados de
la mezcla que antes viajaban \textit{in situ}. Esto no significa que esos puestos vayan a desaparecer,
pero sí que el sector tendrá que adaptarse, y es responsabilidad de quienes adoptan estas
tecnologías que ese cambio no derive en precariedad laboral~\cite{acemoglu2019automation}.

Por otro lado, al exponer puertos de control por red, como el puerto~9999 de voctocore, la
consola de RabbitMQ o la API de telemetría, aumenta el número de puntos de acceso que un atacante podría intentar explotar.
Un acceso no autorizado podría interferir con una retransmisión en directo. Por ello, quien despliega el
sistema debe aplicar autenticación en RabbitMQ, restringir los puertos mediante cortafuegos y no
exponer el sistema a Internet sin protección adicional, siguiendo las recomendaciones del
\ac{NIST} Cybersecurity Framework~\cite{nist_csf2}.

Respecto a la privacidad, la telemetría registra únicamente eventos de operación del mezclador
en ficheros locales, sin recopilar datos personales de operadores ni audiencia. Cuando el sistema
se utiliza para retransmitir imágenes de personas, el operador debe cumplir el \ac{RGPD} y la
\ac{LOPD}, así como respetar los derechos de imagen y de propiedad intelectual sobre el contenido
emitido.}

\vspace{0.15cm}
\noindent\textbf{Aspectos económicos}\\
\normalsize{La ventaja económica principal de Voctomix~2.0 es que hace posible el modelo de
producción remota (\ac{REMI}): el núcleo de procesamiento corre en un servidor centralizado y no hace
falta desplazar a todo el equipo técnico al lugar del evento. Eso reduce directamente los costes
de desplazamiento, alojamiento y logística, que en producciones de varios días o en localizaciones
lejanas suelen ser el gasto más importante.

A eso se suma el ahorro en \textit{hardware}: una instalación multicámara profesional con equipos dedicados
(mezclador Blackmagic ATEM, matrices de vídeo SDI, sistemas de grabación profesional) puede
superar los 5.000~\euro\ en adquisición~\cite{blackmagic_atem}, sin contar mantenimiento ni
formación. Voctomix~2.0 replica esas funcionalidades sobre \textit{hardware} convencional con un coste de
despliegue prácticamente nulo, al ser todo \textit{software} libre.

La arquitectura contenerizada permite además reutilizar la misma instalación para distintos
eventos, reduciendo el coste por producción conforme crece el volumen de retransmisiones.}

\vspace{0.15cm}
\noindent\textbf{Aspectos sociales}\\
\normalsize{El impacto social más claro del proyecto es que acerca la producción audiovisual
profesional a organizaciones que normalmente no pueden permitírsela. Al eliminar la barrera del
\textit{hardware} dedicado, conferencias académicas, jornadas de puertas abiertas, actos culturales y
organismos públicos pueden ofrecer retransmisiones de calidad a sus audiencias. El hecho de
publicar el código bajo licencia libre fomenta además que otros puedan mejorarlo y adaptarlo a
sus propias necesidades.}

\vspace{0.15cm}
\noindent\textbf{Aspectos medioambientales}\\
\normalsize{El modelo REMI reduce los desplazamientos de personal y equipos, y con ello las
emisiones de CO\textsubscript{2} asociadas al transporte. Sustituir \textit{hardware} especializado por
\textit{software} sobre equipos de propósito general también reduce la fabricación de dispositivos
dedicados y alarga la vida útil del equipamiento existente.

El lado negativo es que el servidor que ejecuta voctocore trabaja a alta carga de CPU durante toda
la retransmisión. Las mediciones realizadas en este trabajo muestran un consumo neto, medido con el
sensor RAPL~\cite{khan2018rapl}, que crece de 25,1~W con una cámara activa a 32,2~W con cuatro,
siempre por encima del consumo en reposo del servidor. El procesamiento de vídeo RAW en tiempo real
consume más energía que una solución con \textit{hardware} dedicado de bajo consumo, algo habitual
en infraestructuras de servidor~\cite{iea_datacentres_2023}.

Aun así, el balance global es positivo: el ahorro en emisiones por reducir desplazamientos es
mayor que el aumento de consumo eléctrico del servidor, especialmente en producciones con muchas
personas implicadas o localizaciones geográficamente dispersas.}
\vspace{0.3cm}

% ── A3. Análisis detallado ───────────────────────────────────────────
\Large{\textsc{\textcolor{nar}{A3. Análisis detallado de alguno de los impactos}}}
\vspace{0.2cm}

\normalsize{Se analizan en detalle los dos impactos más relevantes: el económico, por ser la
principal ventaja diferencial del modelo REMI, y el ético, por implicar riesgos que hay que
gestionar activamente.}

\vspace{0.15cm}
\textbf{Impacto económico, reducción de costes en producción remota}

\begin{enumerate}
  \item \textbf{Reducción de personal desplazado} En una producción multicámara convencional,
  todo el equipo técnico viaja al lugar del evento. Con el modelo REMI, el realizador, el técnico
  de audio y el operador de gráficos trabajan desde una sala de control centralizada; solo los
  operadores de cámara necesitan estar presentes \textit{in situ}.

  \item \textbf{Eliminación de \textit{hardware} dedicado} Voctomix~2.0 sustituye mezcladores, matrices
  de vídeo y sistemas de grabación profesional por \textit{software} ejecutado en \textit{hardware} convencional. El
  coste de adquisición pasa de varios miles de euros a prácticamente cero, con la única inversión
  en equipos de red y ordenadores de propósito general.

  \item \textbf{Escalabilidad sin sobrecoste} La arquitectura contenerizada permite levantar el
  sistema completo, incluyendo telemetría y mensajería, con un único comando y sin licencias
  adicionales. La misma instalación sirve para todos los eventos de una organización.
\end{enumerate}

\vspace{0.15cm}
\textbf{Impacto ético, ciberseguridad y responsabilidad profesional}

\begin{enumerate}
  \item \textbf{Superficie de ataque en red} El sistema expone varios puertos TCP: control en
  9999, entradas de cámara en 10000--10005, salidas en 11000--15000, consola RabbitMQ en 15672 y
  API de telemetría en 8080. En una red no controlada, cada puerto accesible es un vector de
  ataque potencial. El ingeniero que despliega el sistema debe aplicar reglas de cortafuegos que
  restrinjan el acceso a los equipos autorizados.

  \item \textbf{Autenticación y control de acceso} El protocolo de control TCP de voctocore no
  incluye autenticación en la capa de aplicación: cualquier cliente que llegue al puerto~9999
  puede enviar comandos al mezclador. En una red local controlada esto es aceptable, pero si hay
  acceso remoto conviene añadir un túnel cifrado (\ac{VPN} o \ac{SSH} port forwarding) para
  proteger el canal de control.

  \item \textbf{Responsabilidad ante el desplazamiento laboral} La adopción de sistemas de
  producción remota debe ir acompañada de medidas de formación y reconversión para el personal
  técnico afectado, de modo que la ganancia en eficiencia no se consiga a costa de empeorar las
  condiciones laborales en el sector audiovisual.
\end{enumerate}

\vspace{0.15cm}
\textbf{Consideraciones sobre privacidad}

El sistema de telemetría registra los eventos de operación del mezclador en ficheros locales y
los publica en un \textit{broker} RabbitMQ interno. El operador debe asegurarse de que: (i)~el acceso al
\textit{broker} está protegido con autenticación; (ii)~los registros no contienen datos personales;
(iii)~el sistema no retransmite contenido protegido por derechos de autor sin la autorización
correspondiente. Al mantener toda la telemetría en infraestructura local bajo control del
operador, cumplir el \acs{RGPD} no requiere modificaciones adicionales.
\vspace{0.3cm}

% ── A4. Conclusiones ────────────────────────────────────────────────
\Large{\textsc{\textcolor{nar}{A4. Conclusiones}}}
\vspace{0.2cm}

\normalsize{Desde el punto de vista ético, social, económico y ambiental, Voctomix~2.0 tiene un
balance global positivo, aunque hay riesgos que deben gestionarse activamente.

En el plano ético, el uso de \textit{software} libre garantiza que el código es transparente y auditable,
y la telemetría local protege la privacidad de los usuarios. Los dos riesgos principales son la
exposición de puertos de red, que exige configurar bien la seguridad, y el posible impacto sobre
el empleo técnico presencial, que obliga a que la adopción de estas herramientas sea responsable.

Económicamente, la mayor ventaja es la reducción de costes de desplazamiento que permite el
modelo REMI, junto con la eliminación del \textit{hardware} dedicado. Esto hace el sistema especialmente
útil para organizaciones con presupuesto limitado.

Socialmente, el proyecto acerca la producción audiovisual profesional a comunidades académicas y
organizaciones que de otro modo no tendrían acceso a ella.

Medioambientalmente, reducir los desplazamientos implica una bajada real en emisiones de
CO\textsubscript{2}, y ese ahorro supera el aumento en consumo eléctrico del servidor. El balance
ambiental es positivo, siendo el ahorro en transporte el factor determinante.

En conjunto, los riesgos identificados (exposición de puertos, impacto sobre el empleo presencial y
mayor consumo eléctrico del servidor) son gestionables con medidas conocidas de seguridad,
formación y eficiencia, y no cuestionan la viabilidad del modelo REMI que propone este proyecto.}
